%% ma: L12-B14-0083
%% lop: 12
%% chuong: 5
%% bai: 14
%% dang: 12.14.5
%% loai: TN4
%% muc: TH
%% dap_an: "A"
%% nguon: "To_n_12_PT_M_t_Ph_ng.pdf (Đề kiểm tra đánh giá kiến thức, Chuyên đề Phương trình mặt phẳng, Nguyễn Anh Khôi - Lê Minh Khôi), Câu 39"
%% kien_thuc: "Bài 14 (phương trình mặt phẳng qua một điểm, biết vectơ pháp tuyến), Bài 7 (hình chiếu của điểm trên trục, mặt phẳng tọa độ)"
%% trang_thai: chinh-thuc
%% ghi_chu: "Tài liệu gốc không kèm đáp án; đáp án do Claude giải và kiểm bằng sympy."
\begin{ex}
Trong không gian $Oxyz$, cho điểm $M(2;0;1)$. Gọi $A$, $B$ lần lượt là hình chiếu của $M$ trên trục $Ox$ và trên mặt phẳng $(Oyz)$. Viết phương trình mặt phẳng trung trực của đoạn $AB$.
\choice{\True $4x-2z-3=0$}{$4x-2y-3=0$}{$4x-2z+3=0$}{$4x+2z+3=0$}
\loigiai{$A(2;0;0)$, $B(0;0;1)$, trung điểm $I\left(1;0;\dfrac12\right)$, $\overrightarrow{AB}=(-2;0;1)$. Mặt phẳng trung trực: $-2(x-1)+\left(z-\dfrac12\right)=0\Leftrightarrow 4x-2z-3=0$. Chọn A.}
\end{ex}
